\section{The Regret solver}
\label{sec:solver}

\subsection{Game representation}

We implement a flop-subgame solver for HUNL. The game tree is built lazily
during CFR traversal: states encode hole cards, board, pot size, stack
sizes, street, and abstracted action history. Terminal nodes are detected
by fold, all-in resolution, or river completion.

\subsection{Abstraction configuration}

\paragraph{Card abstraction.} We implement equity-percentile bucketing.
For each (hole, board) pair, equity is estimated via Monte Carlo sampling
(1{,}000 rollouts by default). Hands are sorted by equity and assigned to
one of $K$ buckets. We evaluate $K \in \{10, 25, 50, 100, 200\}$.

\paragraph{Action abstraction.} We test configurations of $\{2, 3, 5\}$
bet sizes drawn from $\{25\%, 50\%, 75\%, 100\%, \mathrm{all\text{-}in}\}$.

\subsection{Exploitability computation}

We compute exploitability via best-response computation using the
imperfect-information algorithm of \citet{johanson2011efficient}: opponent
reach probabilities are propagated forward, then per-infoset action values
are aggregated and the argmax taken bottom-up. The computation is exact
within the abstracted game and is invoked every $50$ iterations.

\subsection{Parallelization}

Version~1 implements single-threaded alternating updates (player~0 on even
iterations, player~1 on odd). Parallel external-sampling CFR is left for
future work.
